\section{Step 101 Admit a complete picture request}\label{h263-picture-s01}
Check all controls and complete mapped extents before reading current or reference
samples. Check the initial context cursor is zero. An I-picture has no reference
input; reject a nonzero reference field rather than silently choosing a different
interpretation. Retain the caller's output pointers separately from workspace
pointers so nested calls cannot accidentally publish early.

\paragraph{Observe} A high-word mode value whose low bits are one still fails the full-width selector check.
\Needspace{8\baselineskip}
\paragraph{Check} Which input is absent for an I-picture?
\begin{itemize}
  \item[A.] The old reconstructed reference.
  \item[B.] The current raw picture.
  \item[C.] All three current planes.
\end{itemize}

\section{Step 102 Write the picture header into the private context}\label{h263-picture-s02}
Initialize a private writer context over the encoded staging arena at cursor zero.
Emit the 22-bit picture start code, eight-bit temporal reference, thirteen-bit
baseline PTYPE, five-bit quantizer and zero CPM/PEI fields. Select PTYPE from the
admitted picture control. Check every child return even though the derived arena
bound makes capacity failure unexpected for a correct implementation.

\paragraph{Observe} The baseline PTYPE values here are 4160 for I and 4176 for P.
\Needspace{8\baselineskip}
\paragraph{Check} Which context receives the header first?
\begin{itemize}
  \item[A.] The caller context with its count already committed.
  \item[B.] The private workspace context.
  \item[C.] The previous picture reference.
\end{itemize}

\section{Step 103 Begin a row with the right predictor state}\label{h263-picture-s03}
Start row zero immediately after the picture header. For every later row, align
with zero stuffing and emit GBSC, row number, GFID and the unchanged quantizer.
Reset both left-vector components to zero. Do not insert a separate group-zero
header. Keep GFID supplied by the clip controller rather than deriving it from
the row or temporal reference.

\paragraph{Observe} The first macroblock of every explicit row group has predictor zero.
\Needspace{8\baselineskip}
\paragraph{Check} What does an unchanged PTYPE require of GFID?
\begin{itemize}
  \item[A.] A change at every macroblock.
  \item[B.] The row number modulo four.
  \item[C.] The same value as the preceding picture.
\end{itemize}

\section{Step 104 Search the reconstructed reference}\label{h263-picture-s04}
For a P-picture, position current luma at the macroblock origin and pass the
complete old reconstructed Y plane to the earlier bounded search. Use radius two,
current/reference stride 176 and the selected arithmetic mode. Retain the signed
full-pixel result across prediction and block calls. An I-picture skips this
search and uses INTRA blocks.

\paragraph{Observe} A search result of minus one pixel becomes minus two at the syntax boundary.
\Needspace{8\baselineskip}
\paragraph{Check} Which plane supplies reference samples?
\begin{itemize}
  \item[A.] The preceding locally reconstructed Y plane.
  \item[B.] The preceding raw input Y plane.
  \item[C.] The partially written current reconstruction.
\end{itemize}

\section{Step 105 Build six prediction records}\label{h263-picture-s05}
Use complete current and reference Y/Cb/Cr origins and their tight strides in
the six-block predictor request. Pass the selected full-pixel displacement and
current luma origin. Check its return before consuming any prediction record.
For INTRA, pass zero prediction pointer and stride to the block encoder instead
of pretending that an uninitialized prediction record is a zero block.

\paragraph{Observe} Odd whole-luma-pixel motion can require chroma interpolation.
\Needspace{8\baselineskip}
\paragraph{Check} What is the tight chroma-plane stride?
\begin{itemize}
  \item[A.] 176 bytes.
  \item[B.] 88 bytes.
  \item[C.] Six block records.
\end{itemize}

\section{Step 106 Collect levels and position reconstructed pixels}\label{h263-picture-s06}
Visit all six blocks. Derive the current plane position from row, column and
block number, then invoke the complete block encoder. Copy its 128-byte selected
level array into the corresponding macroblock level record. Copy its sixty-four
reconstructed bytes as eight rows into the private destination plane. Advance
each destination row by 176 for Y or 88 for chroma.

\paragraph{Observe} The second Y block starts eight bytes to the right, not sixty-four bytes after the first picture origin.
\Needspace{8\baselineskip}
\paragraph{Check} Which block output is copied into the new picture?
\begin{itemize}
  \item[A.] The selected levels as raw bytes.
  \item[B.] The forward coefficients.
  \item[C.] The clipped reconstructed pixels.
\end{itemize}

\section{Step 107 Emit macroblock syntax and retain actual motion}\label{h263-picture-s07}
Build the syntax request from all six selected level arrays. Use twice the
selected full-pixel displacement and the left predictor in half-sample units.
Permit skip only through the established syntax contract. After success, retain
the actual selected vector for the next column. Never update the left predictor
to the normalized vector difference: those values serve different purposes.

\paragraph{Observe} A skipped macroblock retains actual vector zero even after a nonzero left vector.
\Needspace{8\baselineskip}
\paragraph{Check} Which vector becomes the next left candidate?
\begin{itemize}
  \item[A.] The actual encoded vector.
  \item[B.] The MVD table index.
  \item[C.] The raw unwrapped difference.
\end{itemize}

\section{Step 108 Publish both complete picture results}\label{h263-picture-s08}
After column ten, advance the row; after row eight, append final zero alignment.
Read the complete private byte count and compare it with caller capacity. On
success copy the used stream window and full reconstruction, then commit the
context count. On any earlier failure return zero byte count and retain both
caller arenas. Workspace is deliberately disposable on rejection.

\paragraph{Observe} A capacity-zero call may perform private encoding work but changes neither published output.
\Needspace{8\baselineskip}
\paragraph{Check} When does the caller used count change?
\begin{itemize}
  \item[A.] Immediately after the picture start code.
  \item[B.] After both final output copies succeed.
  \item[C.] Before validating the reference extent.
\end{itemize}

\section{Step 109 Decode and reconstruct the public cases}\label{h263-picture-s09}
Run the 54 public cases on the completed composition. The checker decodes every
header and macroblock, independently repeats the bounded search and computes
floating transform references. It checks the actual selected levels and the
published reconstruction with the established one-unit block accuracy allowance.
It accepts different legal entropy codes and skip choices, private frames and
accurate alternative transforms. It also checks rejection, saved state and guards
in both byte orders.

\paragraph{Observe} The moving-pattern case contains ninety nonzero selected vectors, making a constant-zero motion shortcut observable.
\Needspace{8\baselineskip}
\paragraph{Check} What does the picture checker compare?
\begin{itemize}
  \item[A.] Only the picture start code.
  \item[B.] Only the printed byte count.
  \item[C.] Decoded motion/levels, accurate reconstruction and state effects.
\end{itemize}

\section{Step 110 Keep clip publication separate from picture publication}\label{h263-picture-s10}
The next outer loop reads one complete raw picture, invokes this entry, writes
the encoded bytes and then swaps old/new reference buffers only after accepted
output. The tested full-clip experiment also writes reconstruction planes for
comparison with an independent decoder. A file failure can leave a partial file,
so the loop must stop and report failure rather than treating that file as a
finished clip. Picture memory publication alone does not establish file acceptance.

\paragraph{Observe} All eight proposed pictures encode and independently decode without requiring an all-INTRA substitute.
\Needspace{8\baselineskip}
\paragraph{Check} When may the outer loop swap its reference buffers?
\begin{itemize}
  \item[A.] After the completed picture and required output writes are accepted.
  \item[B.] As soon as the first block is reconstructed.
  \item[C.] Before reading the current raw picture.
\end{itemize}
